%!TEX root = ../Polito PhD thesis template.tex
% ******************************* Chapter 1 ************************************
% NOTA: general introduction to the whole thesis, in three sections: overview
% (the two reactor families and the problem they share), motivations (the
% reference and the representation that stands between it and the many
% evaluations design and operation require), aim and outline. The theory of
% each Part stays in its opening chapter (Chapters 2 and 6). Every statement on
% a chapter is drawn from that chapter; every context statement is cited from
% the bibliographies of the Parts (refs_mg, refs_sph, refs_eg, refs_fm).
% ==============================================================================

\chapter{Introduction}\label{ch:introduction}

\section{Overview}

The growing demand for energy, and the requirement that it be supplied without
carbon emissions, is leading governments and companies to look for solutions
beyond the ones in service today. In the nuclear industry that search has taken
two directions at once. One is the fourth generation of power reactors, which
aims at a better use of the fuel and at a reduced burden of long-lived waste.
The other is a change of scale and of market: reactors of a few megawatts,
transportable, addressed to sites that a large plant has never served, such as
remote communities and installations off the grid, mining operations, ports and
data centres, where the term of comparison is a diesel
generator~\cite{hpmrMarkets}. Westinghouse Electric Company has proposed a
design in each direction, a lead-cooled fast reactor of \SI{950}{MW_{th}}
developed with Argonne National Laboratory~\cite{wecLFR} and the eVinci
micro reactor, in which the heat is removed by alkali-metal heat pipes rather
than by a circulating coolant~\cite{eVinci}.

Designs of this kind have no operating experience behind them, and their
performance has to be predicted, both while the design is being settled and
when it is presented for licensing, by calculation. Experiments remain
indispensable, and are neither quick nor cheap, so that the computing power now
available is used to reduce their number and to decide what they should
measure. The tools that answer these questions, however, were for the most part
built for light-water reactors, and they do not transfer unchanged to a fast
spectrum or to a core made of heat pipes and hydride moderators. This thesis is
concerned with the methods rather than with a single design: it develops
techniques that improve the simulation of concepts of these two kinds, working
on models that are in the public domain and stating in each case how far the
result carries beyond the core on which it was obtained.

The reactors studied here therefore belong to two families that have, at first
sight, little in common. Lead-cooled fast reactors (LFRs) operate at low
pressure and at temperatures higher than those of light-water reactors, and
their fast spectrum fissions plutonium and the minor actinides, so that the same
core produces energy and reduces the long-lived waste. Part~\ref{part:lfr}
considers two of them, a simplified model of the Westinghouse
LFR~\cite{wecLFR} and the ALFRED core, the European demonstrator of the
technology, in the specification of the benchmark coordinated by the Nuclear
Energy Agency~\cite{neaBench}. Heat-pipe micro
reactors (HPMRs), the subject of Part~\ref{part:hpmr}, are cores of a few
megawatts in which heat pipes take the place of the circulating coolant and the
reactivity is governed by drums rotating in the radial reflector. Since the
commercial designs of this class are not published in the detail that a
neutronic study requires, the reactor considered is the micro reactor developed
at Argonne National Laboratory as a modelling exercise meant to collect in one
core the features that characterize the class, heat pipes, TRISO fuel,
yttrium-hydride moderator pins and control drums~\cite{Stauff2024}, whose models
are distributed on the Virtual Test Bed~\cite{moose_vtb_mrad_description}.

What the two families share is a problem. The computational chain that decades
of light-water reactor practice have consolidated proceeds from the point-wise
nuclear data to a few-group set of homogenized constants~\cite{cacuci2010handbook}
and from there to a coarse-mesh nodal solution of the whole
core~\cite{Lawrence1986,Smith1986}. On it rest the core simulators with which
those plants are designed and the online monitoring systems with which they are
operated~\cite{BoydMiller1996}, and its accuracy rests on an equivalence between
the heterogeneous assembly and its homogenized representation that was
established for light-water lattices~\cite{Smith1986}. Neither family inherits
that chain as it stands. In a fast system the choices that light-water practice
settled long ago, the environment in which the weighting flux is computed and
the mesh over which the energy is condensed, no longer follow from that practice
and have to be examined again; and the computing power now available makes it
possible to examine them on the design itself, with the group constants
generated by a Monte Carlo calculation of the core rather than by the chain of
lattice steps, which is the route taken in Part~\ref{part:lfr}. For a
micro reactor the developers of the reactor physics tools of the NEAMS program
observe that heat pipes, hydride moderators, TRISO fuel and control drums call
for a geometrical generality that the traditional nodal methods do not
possess~\cite{Wang2025griffin,Stauff2024}, that the homogenized solution departs
from the heterogeneous one more than it does in light-water
reactors~\cite{Pandya2022}, and that the cross sections have to be tabulated
against the angle of the drums~\cite{Roskoff2022physor}. With no operating
experience to appeal to, the reference against which a model of either family is
judged is the explicit Monte Carlo solution of the full core.

\section{Motivations}

The Monte Carlo method occupies a place apart among the tools of reactor
physics. It solves the transport problem by simulating a large number of neutron
histories in the true geometry of the system, with the nuclear data at their
tabulated energies rather than averaged over energy groups, so that what limits
its accuracy is the number of histories and the data themselves, and not a
simplification of the physics. Its general-purpose
implementations~\cite{mcnp63,Serpent_code_description,openmc} are for this
reason the standard against which every approximate description of a core is
measured, time-dependent modes included~\cite{serpentDynamic}.
That the accuracy of these codes is
paid in computing time is well known, but the form the cost takes for the
problems of this thesis is worth stating. Neither the design nor the operation
of a reactor asks for one state of the core. A design space is explored over
the positions of the control system, the temperatures and the states of the
fuel; a fuel cycle is a sequence of states; the monitoring of an operating core
requires its critical configuration to be predicted continually as the core
evolves. For the micro reactor of Part~\ref{part:hpmr} the Monte Carlo critical
state of one configuration costs several tens of CPU hours, an expense that a
study of a few states absorbs and a study of thousands does not, and where the
answer is needed within the time of operation the size of the machine is not
the limiting factor.

Every practical scheme therefore interposes, between the reference and the many
evaluations, a representation of the reactor that is built once from a few
reference calculations and consulted many times. In Part~\ref{part:lfr} it is
the set of multi-group constants: the cross sections condensed over a few
energy groups and homogenized over the regions of the core, on which a
deterministic solver runs in a fraction of the time. In Part~\ref{part:hpmr} it
is a database of fission matrices: a Monte Carlo calculation tallies, at each
of a set of reference states, the matrix that couples every fissile region of
the core to every other, and any intermediate state is then recovered by
interpolating the matrices and solving an eigenvalue problem of modest
size~\cite{CarneyANE2014,WaltersNSE2018}. The two representations have nothing
in common in what they store or in the solver that uses them, but they share
what decides their usefulness. Each is obtained from the reference by an
operation, a weighting over energy and space in one case, a tally at fixed
conditions in the other, that discards information the reference possessed,
and the calculations downstream cannot restore it: what is lost in building
the representation is lost for every evaluation that follows, and no accuracy
of the solver recovers it.

The thesis works on the representation itself, along two lines that recur in
both Parts. The first is the choice of how it is built. Part~\ref{part:lfr}
asks by which route the constants should be generated, from one full-core
model or from a set of mini-core models, and on which energy mesh they should
be condensed; Part~\ref{part:hpmr} asks which parameters the database should
span, how dense it has to be and how its matrices should be interpolated. The
second line is the correction: a quantity computed from a small number of
reference calculations and applied to the representation so that it reproduces,
at the conditions of the reference, what the reference reproduces. In
Part~\ref{part:lfr} these are the superhomogenization factors~\cite{Hbert_1980},
which rescale the constants of a region until the homogenized calculation
returns the reaction rates of the heterogeneous one. In Part~\ref{part:hpmr}
they are the correction ratios, ratios taken element by element between two
matrices tallied at reference conditions that differ in a single effect, by
which the interpolated matrix is multiplied to account for the burnup, for the
temperature of the reflector and for the temperature gradients inside the core.
Both act on the data and not on the solver, which receives corrected input and
requires no modification, and this is what makes them applicable to a code that
cannot be modified, because it is qualified or because its source is not
available.

Two rules govern the use of the representation throughout. The first is that
the reference keeps the last word. No result of this thesis is a prediction of
the intermediate model alone: every energy structure, every critical state and
every corrected matrix is set against a Monte Carlo calculation made for that
purpose, and the rule holds also where a further layer of approximation is
introduced for the sake of speed. In Chapter~\ref{ch:energy-grid} a neural
network trained on the outputs of the deterministic solver takes the place of
the solver inside an optimization loop; the network is allowed to explore, and
the solver alone decides, since every structure the search proposes is
recomputed with it.

The second rule concerns the nature of that reference. A Monte Carlo result is
not an exact number but an estimate carrying a statistical uncertainty, and a
comparison that ignores it credits a model with an agreement that the reference
itself cannot resolve. Two energy structures whose predictions differ by less
than the uncertainty of the reference are not, on that evidence, one better than
the other, and it is for this reason that Chapter~\ref{ch:energy-grid} fixes the
thresholds of its objective function in relation to that uncertainty. The same
applies downstream: the coefficients of a fission matrix are tallied quantities
and carry their own noise, which the interpolation passes on to every result
obtained from them, so that Part~\ref{part:hpmr} derives an estimate of the
uncertainty thus inherited and checks it against independent repetitions of the
whole database.

\section{Aim of the work and outline of the thesis}

The aim of this thesis is to improve the accuracy, and to contain the cost, of
the neutronic analysis of two advanced reactor concepts by acting on the
representation that stands between the Monte Carlo reference and the
calculations that their design and operation require. For the lead-cooled fast
reactor that representation is the multi-group data of a deterministic solver,
and the work concerns the route by which the constants are generated, their
correction and the energy mesh on which they are condensed. For the heat-pipe
micro reactor it is a database of fission matrices, and the work concerns the
construction of the database, its use for the search of the critical state and
the corrections that extend it to the conditions of an operating core. The
thesis is organized in two Parts, each opened by a chapter that provides the
theory common to the works it contains, and closed by a chapter of conclusions
common to both.

Chapter~\ref{ch:part1-opener} opens Part~\ref{part:lfr} and lays the common
ground of its two works: the passage from the transport equation to its
multi-group form, the definition of the group constants as flux-weighted
averages over an energy interval and over a region of the core, the assumptions
that this definition requires and the errors it introduces. Two of the degrees
of freedom that the definition leaves open, the environment in which the
weighting flux is computed and the mesh over which the energy is condensed, are
the subjects of the two chapters that follow. In both, the multi-group
constants and the reference solution come from the same Serpent~2 calculations,
so that what is measured is the multi-group approximation and not the nuclear
data.

Chapter~\ref{ch:sph-lfr} addresses the Westinghouse LFR. Its constants are
generated along two routes, from a single full-core model and from a set of
mini-core models each representing one region of the system, and they feed a
diffusion solver written in FreeFEM++ on a finite-element mesh.
Superhomogenization factors are then computed for the regions responsible for
the largest discrepancies in the power distribution, and the two routes, with
and without correction, are compared on the multiplication factor, the power
distribution and the kinetic parameters, and on the computational time each
requires when several configurations of the control rods have to be analyzed.

Chapter~\ref{ch:energy-grid} addresses the ALFRED core and the energy mesh
itself. The boundaries of the few-group structure become the unknowns of an
optimization: a genetic algorithm searches the subsets of the boundaries of a
thirty-group fine mesh for the structure whose deterministic solution
reproduces, in three configurations of the control rods at once, the
multiplication factor, the delayed neutron fraction, the neutron lifetime, the
power distribution and the flux distribution of the reference. The work that
formulated the problem had run three searches from three random seeds and
obtained three different structures; this chapter runs one thousand. That
number is within reach because a neural network trained on $91\,723$ full-core
evaluations of the deterministic code FRENETIC~\cite{frenetic,freneticBench}
stands in for the code inside the search, while every structure the searches
return is recomputed with the code. The chapter measures how strongly the
outcome depends on the seed and how many searches are needed, identifies the
boundaries that the optimized structures have in common, compares the best of
them with the structures obtained with the solver inside the loop and with
structures built by standard spacing rules, and states the limits within which
the result holds.

Chapter~\ref{ch:fm-theory} opens Part~\ref{part:hpmr}. It derives the fission
matrix from the transport equation, reviews the strategies by which the matrix
can be generated, describes the reactor model and sets up the framework that the
two following chapters share: the parametrization of the database over the
operating state, the treatment of the statistical uncertainty of the
interpolated results and the metrics by which the errors are measured.

Chapter~\ref{ch:crit-search} applies that framework to the criticality search:
finding, at a given power and at a given point of the fuel cycle, the angle of
the control drums that makes the core critical, together with the fission
source and the temperature field that correspond to it. A database of fission
matrices is built over the fuel temperature and the drum angle, its density is
fixed by a convergence study, and a correlation between power and fuel
temperature stands in for the thermal solver, so that the search runs on the
fission matrix alone. The algorithm is verified against reference Monte Carlo
calculations at beginning of life and then carried along the fuel cycle by
means of the Burnup Correction Ratio; the assumptions on which that extension
rests are examined one by one, together with the density of the snapshot grid
and the computational cost of the whole procedure.

Chapter~\ref{ch:correction-ratios} removes the restriction to isothermal
states. The database of Chapter~\ref{ch:crit-search} has every material at one
temperature, and two effects of an operating core lie outside its reach, a
reflector colder than the core and a temperature that varies sharply from one
cell to the next. Each receives a correction ratio built on the scheme of the
burnup correction, with a weight that scales the ratio by the fraction of the
reference effect present at the target state: the Reflector Correction Ratio,
developed on the three-dimensional model of the criticality search, and the
Core Temperature Correction Ratio, developed on a two-dimensional variant of
the same design. The chapter also carries out the analysis of the statistical
uncertainty of the interpolated matrix announced in Chapter~\ref{ch:fm-theory}:
the correlation among the tallied coefficients is modelled, its invariance with
temperature is examined and the resulting estimate is validated against
independent replicas of the database, before the results of the core
temperature correction are read against the combined uncertainty.

Chapter~\ref{ch:conclusions} gathers the conclusions of the two Parts and the
developments that the work leaves open.
